\section*{Chapitre \ref{ch:s2:power-and-gas-trading} --- Négociation d'électricité et de gaz}

\begin{solution}{exo:s2:power-and-gas-trading:1}
$1.5 \times 3.05 = 4.575$~: environ 4.6~\euro par MWh sous le prix day-ahead, avant les autres bruits.
\end{solution}

\begin{solution}{exo:s2:power-and-gas-trading:2}
Charger 2~MWh de stockage à 0~\euro ne coûte rien~; décharger restitue $2\sqrt{0.88} = 1.876$~MWh au réseau (le côté charge a perdu l'autre
facteur $\sqrt{0.88}$, payé à 0~\euro), vendu à 100~\euro~: 187.62~\euro.
\end{solution}

\begin{solution}{exo:s2:power-and-gas-trading:3}
$2.54 \times 0.61 \times 24 \times 727 / 2 \approx 13\,500$~: environ 13\,500~\euro par an~; avec les données non arrondies, 13\,600~\euro.
\end{solution}

\begin{solution}{exo:s2:power-and-gas-trading:4}
L'edge est d'environ 3.5~\euro par MWh avant coûts (2.54~\euro après une demi-fourchette de 1~\euro)~; une demi-fourchette de 3~\euro en absorbe l'essentiel, il reste 0.54~\euro.
Un edge de prévision mesuré en quelques euros est faible face à des coûts de quelques euros.
\end{solution}

\begin{solution}{exo:s2:power-and-gas-trading:5}
La production solaire culmine à midi~; lorsque le solaire et l'éolien dépassent ensemble la demande augmentée de la capacité d'exportation, les producteurs qui doivent fonctionner ou fonctionneront quoi qu'il arrive (production subventionnée,
centrales inflexibles) soumettent des offres sous zéro~. En 2024--2025, 23.4~\% des heures à 13:00 étaient négatives, aucune à 19:00.
\end{solution}

\begin{solution}{exo:s2:power-and-gas-trading:6}
Ils posaient des questions différentes, avec des données et des modèles différents~: Kiesel et Paraschiv font dépendre la réponse d'un seuil portant sur la demande attendue
couverte par la capacité conventionnelle, de sorte que l'asymétrie n'apparaisse que dans certains états~; Ziel a testé si les erreurs positives et négatives diffèrent dans son
modèle et n'a trouvé aucune différence significative. Un effet qui n'existe qu'au-dessus d'un seuil peut disparaître dans une moyenne.
\end{solution}

\begin{solution}{exo:s2:power-and-gas-trading:7}
Le revenu day-ahead passe de 69\,900~\euro à 83\,800~\euro, soit 13\,800~\euro de plus, et le supplément infrajournalier passe de 2\,900~\euro à 4\,900~\euro. Le tableau ne tient pas compte
de la dégradation~: chaque cycle supplémentaire use la durée de vie de la batterie, et le second cycle de la journée utilise des écarts de prix plus faibles.
\end{solution}

\begin{solution}{exo:s2:power-and-gas-trading:8}
Le backtest planifiait chaque jour en connaissant les prix d'enchère de ce jour-là, que doit prévoir un opérateur réel avant de soumettre ses offres~; il n'intégrait ni dégradation, ni
fourchette, ni frais de réseau. Il s'agit d'une borne supérieure pour le seul marché day-ahead, et non d'une prévision~; les revenus infrajournaliers et de réserve, que le calcul omet,
poussent en sens inverse.
\end{solution}

\begin{solution}{pb:s2:power-and-gas-trading:1}
\begin{enumerate}
\item Le prix day-ahead moins le prix infrajournalier pour la même heure, déterminé par ce qui a changé entre l'enchère et la livraison.
\item 3.05~\euro par MWh et par GW de charge résiduelle, corrélation 0.83, sur 2024--2025.
\item Le prix day-ahead moins 3.05~\euro par GW de vent inattendu, erreurs de 2~GW d'écart-type avec une autocorrélation horaire de 0.9, plus 6~\euro de
bruit~; le tout est introduit volontairement.
\item Kiesel et Paraschiv~: les prix infrajournaliers s'ajustent de façon asymétrique aux erreurs de prévision sur les énergies renouvelables et au volume, selon un seuil de demande. Ziel~: aucune
différence significative entre les effets des erreurs positives et négatives.
\item Vendre en day-ahead et racheter en infrajournalier lorsqu'on anticipe plus de vent que n'en suppose l'enchère~; l'inverse lorsqu'il y en a moins~; aucune position sur les petites prévisions.
\item Compétence 0.5, demi-fourchette 1~\euro~: 2.54~\euro par MWh, 13\,600~\euro par an et par MW, IC 0.36~; compétence 0.3~: 1.08~\euro, 5\,800~\euro, 0.21~; demi-fourchette
3~\euro~: 0.54~\euro, 2\,900~\euro.
\item Une \omterm{def:s2:power-and-gas-trading:weather}{position météo} sur l'électricité ou le gaz, fondée sur une prévision qui diverge de celle du marché, clôturée lorsque les prévisions convergent ou que la météo se réalise.
\item De meilleures prévisions publiques et des coûts moindres.
\item Planifier une batterie pour charger quand l'électricité est bon marché et décharger quand elle est chère, en day-ahead puis en infrajournalier, dans ses limites.
\item Bon marché à midi, cher le soir~; 457 heures négatives en 2024 et 573 en 2025, soit 23.4~\% des heures à 13:00.
\item Programmation dynamique quotidienne sur l'état de charge et l'énergie chargée, un cycle, rendement aller-retour de 88~\%, batterie vide en fin de période.
\item Day-ahead 69\,900~\euro par an~; supplément infrajournalier 2\,900~\euro (3\,800~\euro sans fourchette, 1\,600~\euro à 3~\euro)~; deux cycles 83\,800~\euro et
4\,900~\euro.
\item 69\,900~\euro par MW et par an en day-ahead, 72\,800~\euro avec réoptimisation infrajournalière (2\,900~\euro de plus), sur les prix allemands de 2024--2025.
\item La planification day-ahead connaît les prix de l'enchère~; la planification infrajournalière connaît d'emblée les prix infrajournaliers de la journée.
\item Les erreurs introduites volontairement déplacent les prix de quelques euros, alors que l'amplitude quotidienne moyenne dépasse la centaine d'euros.
\item Erreurs de prévision corrélées d'une heure à l'autre, sauts de prix, pénalités de déséquilibre, limites de l'actif et dégradation.
\item Soumettre des offres sur des prévisions de prix établies avant l'enchère, réoptimiser en infrajournalier à mesure que les prix arrivent, et intégrer la dégradation et les frais dans le calcul.
\item \omterm{def:s2:power-and-gas-trading:battery}{Arbitrage par batteries} en day-ahead~: davantage de batteries aplatissent la forme quotidienne qu'il négocie.
\item Les deux sont rémunérés par les acteurs contraints de négocier (flux au fixing, producteurs qui compensent leurs erreurs), et les deux exigent une prévision de ce besoin.
\item La livraison à une heure donnée, et l'information sur ce dont cette heure aura besoin.
\end{enumerate}
\end{solution}

\begin{solution}{iq:s2:power-and-gas-trading:1}
Certains producteurs sont payés pour produire (subventions par MWh) ou ne peuvent pas s'arrêter à moindre coût (centrales inflexibles), si bien qu'ils soumettent des offres négatives plutôt que de réduire leur production~; lorsque
leur offre dépasse la demande et la capacité d'exportation, l'enchère se conclut sous zéro.
\end{solution}

\begin{solution}{iq:s2:power-and-gas-trading:2}
Prévoir la variation du vent, du solaire et de la charge entre les prévisions de l'enchère et les siennes, la traduire en prix par la pente de la charge résiduelle dans
l'état de l'ordre de préséance économique (merit order) de cette heure, et la confronter aux spreads réalisés~; intégrer les indisponibilités de centrales et les flux d'interconnexion.
\end{solution}

\begin{solution}{iq:s2:power-and-gas-trading:3}
Vendre ces heures en day-ahead et prévoir de racheter en infrajournalier, dimensionné d'après le bilan de la prévision à cet horizon et l'effet attendu sur les prix
(environ 3~\euro par GW ici), et dans des limites tenant compte des erreurs corrélées entre les heures de l'après-midi.
\end{solution}

\begin{solution}{iq:s2:power-and-gas-trading:4}
Limites sur la position ouverte par heure et par bloc, sur la somme des heures corrélées, sur l'exposition au déséquilibre à l'heure de fermeture du guichet, sur la perte quotidienne, et les tests de résistance
pour les pics de prix et les défaillances de prévision.
\end{solution}

\begin{solution}{iq:s2:power-and-gas-trading:5}
Un service de prévision des prix~; un optimiseur par batterie (programmation dynamique ou programme linéaire) relancé à chaque nouveau prix~; un gestionnaire des positions qui
compense le parc par rapport aux marchés~; un routage des ordres vers le marché infrajournalier~; une télémesure de l'état de charge et de la disponibilité~; et un planning de secours.
\end{solution}

\begin{solution}{iq:s2:power-and-gas-trading:6}
Acheter 1~MWh à $p_1$ restitue $\eta$~MWh vendus à $p_2$~: profit $\eta p_2 - p_1 > 0$ si et seulement si $p_2 > p_1/\eta$. Avec $p_1 < 0$, la batterie est payée pour
charger, si bien que l'opération est rentable dès que $\eta p_2 > p_1$, y compris pour certains $p_2$ négatifs~; la contrainte porte alors sur la capacité, et non sur le prix.
\end{solution}
